\documentclass[12pt,a4paper]{article}
\usepackage[utf8]{inputenc}
\usepackage[T1]{fontenc}
\usepackage{amsmath}
\usepackage[left=2.00cm, right=2.00cm, top=2.00cm, bottom=2.00cm]{geometry}
\usepackage{amssymb}
\usepackage[italian]{babel}

\title{Tutorato 3 - Analisi\\Complementi di Matematica per le Scienze Chimiche\\ A.A. 2025/2026}
\author{Massimiliano Ghiotto - massimiliano.ghiotto01@universitadipavia.it}
\date{}

\usepackage{amsthm}
\theoremstyle{definition}
\newtheorem{exe}{Esercizio}

\begin{document}
\maketitle


% \begin{exe}
%     Verificate che la funzione $f(x,y) = \alpha x + \beta y + \gamma$ è armonica per ogni $\alpha, \beta, \gamma \in \mathbb{R}$.
% \end{exe}

% \begin{exe}
%     Verificate che
%     \begin{itemize}
%         \item $e^{kx} \cos(ky)$ è armonica per ogni $k \in \mathbb{R}$;
%         \item $e^{kx} \cos(ky) + e^{hx} \sin(hy)$ è armonica per ogni $h, k \in \mathbb{R}$;
%         \item $\ln(\sqrt{x^2+y^2})$ è armonica su $\mathbb{R}^2 \setminus \{(0,0)\}$.
%     \end{itemize}
% \end{exe}

% \begin{exe}
%     Dimostrate il teorema enunciato nella teoria: se $f,g : A \subseteq \mathbb{R}^n \to \mathbb{R}$ sono armoniche, allora per ogni $\alpha, \beta \in \mathbb{R}$, $\alpha f + \beta g$ è armonica.
% \end{exe}

\begin{exe}
    Calcolate il laplaciano delle seguenti funzioni scritte in coordinate polari:
    \begin{enumerate}
        \item $g(\rho, \theta) = \rho \theta^2$;
        \item $g(\rho, \theta) = \rho^2 \theta$;
        \item $g(\rho, \theta) = \rho \sin(\theta)$;
        \item $g(\rho, \theta) = \sin(\rho \theta)$;
        \item $g(\rho, \theta) = e^\rho \cos(\rho)$;
        \item $g(\rho, \theta) = f(\rho)$.
    \end{enumerate}
\end{exe}

\begin{exe}
    Determina in quali punti di $\mathbb{R}^3$ la seguente funzione è derivabile, calcolando esplicitamente le derivate in tale caso, e differenziabile. Determina il più grande insieme aperto in cui $f \in C^1$:
    \[
        f(x,y,z) = zx^2 - xy^3 + z^4
    \]
\end{exe}

\begin{exe}
    Determina in quali punti di $\mathbb{R}^2$ la seguente funzione è derivabile, calcolando esplicitamente le derivate parziali in tale caso, e differenziabile. Determina il più grande insieme aperto in cui $f \in C^1$:
    \[
        f(x,y) = |x|y
    \]
\end{exe}

\begin{exe}
    Determina in quali punti di $\mathbb{R}^2$ la seguente funzione è derivabile, calcolando esplicitamente le derivate parziali nell'origine, e differenziabile:
    \[
        f(x,y) = \begin{cases}
            \frac{x^3+2y^4}{x^2+y^2} & \text{per } (x,y) \neq (0,0) \\
            0                        & \text{per } (x,y) = (0,0)
        \end{cases}
    \]
\end{exe}

\begin{exe}
    Data la funzione
    \[
        f(x,y) = \begin{cases}
            \frac{2x^3+x^2y^2+2xy^2+4y^5}{x^2+y^2} & \text{per } (x,y) \neq (0,0) \\
            0                                      & \text{per } (x,y) = (0,0)
        \end{cases}
    \]
    calcola le derivate parziali di $f$ in $(0,0)$ e stabilisci se $f$ è differenziabile nell'origine. Calcola inoltre $D_v f(0,0)$ per $v = \left(\frac{\sqrt{3}}{2}, \frac{1}{2}\right)$.
\end{exe}

\begin{exe}
    Scrivi una forma differenziale esatta e una non esatta prima in $\mathbb{R}^2$, poi in $\mathbb{R}^3$ e infine in $\mathbb{R}^4$.
\end{exe}

\begin{exe}
    Determina se le seguenti forme differenziali sono esatte e, in caso affermativo, determinane un potenziale.
    \begin{enumerate}
        \item $yzw dx + xzw dy + xyw dz + xyz dw$ su $\mathbb{R}^4$;
        \item $(\cos(a)b \cos(c) + c^2) da + \sin(a) \cos(c) db + (-\sin(a)b \sin(c) + 2ac) dc$ su $\mathbb{R}^3$;
        \item $\cos(u) du + \sin(v) dv$ su $\mathbb{R}^2$;
        \item $v \cos(uv) du + u \cos(uv) dv$;
        \item $c |(a,b,c)| da + a |(a,b,c)| db + b |(a,b,c)| dc$;
        \item $x |(x,y,z)| dx + y |(x,y,z)| dy + z |(x,y,z)| dz$.
    \end{enumerate}
\end{exe}

\begin{exe}
    Considera la forma differenziale $|x| dx + |y| dy$ definita su $\mathbb{R}^2$.
    \begin{enumerate}
        \item Dimostra che la forma differenziale è esatta.
        \item Determina qual è l'insieme più grande $D$ in cui ha senso chiedersi se la forma differenziale è chiusa.
        \item Giustifica in base alla teoria e a quanto determinato al punto precedente che la restrizione della forma differenziale all'insieme $D$ è chiusa.
    \end{enumerate}
\end{exe}

\begin{exe}
    Se $U$ è un potenziale del campo $F : \mathbb{R}^3 \to \mathbb{R}^3$ e se $V$ è un potenziale del campo $G : \mathbb{R}^3 \to \mathbb{R}^3$, determina
    \begin{enumerate}
        \item se $U+V$ è il potenziale di $F+G$;
        \item determina quale campo ha come potenziale $U \cdot V$.
    \end{enumerate}
\end{exe}

\begin{exe}
    Determina per quale valore di $k \in \mathbb{R}$ il campo $F(x,y,z) = (3x^2 y^2 z, 2x^3 y z^k, x^3 y^2)$ è conservativo, poi calcola il potenziale corrispondente.
\end{exe}



\end{document}
